\documentclass[a4paper]{slides}
\usepackage[francais,english]{babel}

\newcommand{\mytitle}[1]{\begin{center}\begin{large}#1\end{large}\end{center}}

\sloppy
\selectlanguage{francais}

\begin{document}

\begin{slide}
  \mytitle{TH\`EMES}
  Je presenterais les r\`egles choisis pour illustrer mon expos\'e.\\
  Ces r\`egles sont des automates cellulaires.

  Je presenterais les moteurs \`a tables de transitions que j'ais r\'ealis\'es pour
  calculer des trajectoires dans l'espaces des configurations de
  familles de r\`egles.
  
  Je commenterais la r\'ealisation d'une ex\-p\'e\-ri\-ence de calcul de trajectoire
  et je pr\'esenterais quelques visualisation exp\'erimentales de
  trajectoires.
  
  Enfin, j'aborderais les espaces de r\`egles
  en introduisant un nouveau param\`etre de controle de l'espaces des
  r\`egles\\
  et en montrant la r\'ealisation d'une transformation d'architecture
  SIMD en r\`egle utilisable par un moteur cellulaire.
\end{slide}

\begin{slide}
  \mytitle{R\`EGLES}
  j'ais choisis 3 r\`egles additives \`a voisinnage de MOORE.
  
  Les r\`egles additives peuvent \`etre d\'efinis par une table secondaires
  dont les entr\'ees sont fonctions de la somme des voisins.
  
  Life et Anneal sont de r\`egles binaires.
  Voici leur table secondaire.
  Le U signifie que l'\'etat de la cellule reste unchang\'e.
  
  Les trajectoires pr\'esent\'ees dans l'expos\'e sont obtenu \`a partir d'un
  configuration initiale al\'eatoire.
  
  Qmean est une r\`egle \`a 8 bits par cellule que j'ais d\'ecouvert en
  testant un moteur \`a voisinnage additif virtuel.
\end{slide}

\begin{slide}
  \mytitle{MOTEURS}
  J'ais choisis d'impl\'ementer les moteurs de calcul par tables de
  transitions\\
  et collecte pipelin\'e de l'index de voisinage.
  
  Ce choix permet la r\'ealisation de moteurs rapide sur machine
  s\'equentielle.
  
  Il permet \'egalement la manipulation des r\`egles afin de Mesures et
  controle.\\
  Ce qui permet l'utilisation des moteurs pour l'exploration des espaces
  de r\`egles.
  
  L'inconv\'enient de tables de transitions est leur taille.\\
  Je propose des m\'ethodes de compositions de tables et de voisinage
  virtuels.
  
  Ici le pipeline calculant le voisinage virtuel SUM.
\end{slide}

\begin{slide}
  \mytitle{EXPLORATIONS}
  \begin{footnotesize}
    Cette exp\'erience \`a pour but la mise en \'evidence d'un invariant
    d'echelle espace-temps sur les trajectoires de Anneal.
    
    j'ais calcul\'e des trajectoire de Anneal pour des configuration initiales de
    tailles croissantes.
    
    j'ais mesur\'e la population (le nombre de site \`a 1) et la taille des
    fronti\`eres pour chaque points de trajectoires.
    
    j'ais normalis\'e les mesures.
    
    Le r\'esultat le plus important est la mise en \'evidence de l'existence
    de cycle limite dont la taille est de l'ordre de grandeur de l'ar\`ete
    de la configurations.
    
    Cette figure me permet d'introduire les visualisation:\
    Dans un contexte de visualisation une trajectoire est un espace-temps.
    
    Une extraction d'espace-temps qui pr\'eserve l'espace, mais pas le
    temps.
    
    Une s\'equence d'accumulation sur l'espace qui preserve le temps, mais
    pas l'espace.
  \end{footnotesize}
\end{slide}

\begin{slide}
  \mytitle{VISUALISATION 1 OF 3}

  Je propose de caract\'eriser les visualisations par une combinatoire
  de trois familles d'op\'erateurs:~les extractions, les diff\'erences et
  les accumulations.
  
  j'ais utilis\'e l'op\'erateur xor pour r\'ealiser une diff\'erence entre
  points deux points de temps successif.
  
  Une premi\`ere application permet la suppresion des points fixes
  locaux.
  
  Une deuxi\`eme application la suppression des cycle limite de locaux de
  tailles 2.
  
  l'accumulation additive sur le temps r\'ev\`ele les densit\'es d'occupation
  de l'epace-temps.
\end{slide}

\begin{slide}
  \mytitle{VISUALISATION 2 OF 3}
  
  Voici une autre visualisation possible de l'espace-temps et de
  l'espace-temps apr\`es suppression des points fixes locaux et cycles
  limites locaux de taille 2.
  
  Cette visualisation r\'ealise une extraction par projection.
  
  C'est un visualisation de fin de chaine, elle ne permet pas d'autre
  mesure ou transformations.
\end{slide}

\begin{slide}
  \mytitle{VISUALISATION 3 OF 3}
  Cette page pr\'esente quelques exemples de visualisation.
  
  j'ais r\'ealis\'e des outils de visualisation adapt\'e \`a la nature
  binaires des espace-temps cellulaires\\ et je propose une m\'ethode de
  visualisation par projection relativiste qui associe une variation
  du temps \`a diff\'erents point d'espaces visualis\'e.

  Ici l'augmentation de la distance au centre de la projection
  rapproce de la configuration initiale de la trajectoire.
\end{slide}

\begin{slide}
  \mytitle{ESPACES DE R\`EGLES 1 of 2}
  L'exploration des espaces de r\`egles necessite l'existence d'une
  distance entre les r\`egles qui d\'epends de mesure appliqu\'e sur les
  tables de transitions.
  
  La mesure standard est lamda, qui donne la densit\'e d'une table.
  
  Je d\'efinis flip comme le nombre de changement d'\'etats rencontr\'e en
  balayant une table de transition.\\
  Sa valeur d\'epends de l'ordre d'ecriture de l'index de voisinnage.
  
  je d\'efini le param\`etre d'enchassement comme l'influence de chaque
  voisin sur une r\`egle.
  
  Le calcul est fait pour chaque bit de voisinage par l'application de
  flip sur la s\'equence extraites de la table de transition en fonction
  du poid de ce bit d'index.
  
  Une valeur unique est obtenu par la moyenne de l'enchassement de
  chaque voisin.
\end{slide}

\begin{slide}
  \mytitle{ESPACES DE R\`EGLES 2 of 2}
  J'ais calcul\'e la valeurs de lambda et de l'enchassement pour
  l'ensemble des r\`egles additives \`a 3 fonctions secondaires.
  
  Le pic montre la position de Anneal.
\end{slide}

\begin{slide}
  \mytitle{SIMD}
  J'ais r\'ealis\'e une table de transition pour le PE GAPP.
  
  Le moteur fournit un voisinage virtuel pour la m\'emoire et les
  instructions.
  
  La table de transition \`a une taille de 1 m\'egabyte.
\end{slide}

\begin{slide}
  \mytitle{CONCLUSION}
  En conclusion je propose deux objectifs reposant sur les r\'ealisations
  et proposition que j'ais pr\'esent\'ees:
  
  La r\'ealisation d'une carte des trajectoires des r\`egles
  additives binaires param\'etr\'ee par lambda et l'enchassement.
  
  L'utilisation de tables de transition d'architectures SIMD pour le
  s\'equencement et la combinaison de tables de calculs de trajectoire et
  de visualisation.
\end{slide}

\end{document}

% Local Variables: 
% TeX-master: t
% End: 
